\section{Non-Perturbative Scale Setting Strategy}
\label{sec:scale-setting}

This section describes the standard strategy for defining the physical mass scale in the continuum limit.

\subsection{The Scale Setting Problem}
The classical Yang-Mills Lagrangian contains no dimensionful parameters. The physical theory requires a mass scale (dimensional transmutation).

\begin{definition}[Non-Perturbative Scale Setting]
We define the physical lattice spacing $a(\beta)$ implicitly through a reference physical quantity, specifically the string tension $\sigma$. 
\[
\sigma_{lat}(\beta) = \sigma_{phys} \cdot a^2(\beta)
\]
\end{definition}

\subsection{Requirements for Success}
For this definition to be valid and lead to a non-trivial continuum theory, two conditions must be met:

\begin{enumerate}
    \item \textbf{Confinement:} $\sigma_{lat}(\beta) > 0$ for all $\beta$. This is known at strong coupling (small $\beta$) but is open at weak coupling.
    \item \textbf{Asymptotic Scaling:} As $\beta \to \infty$, the derived $a(\beta)$ must satisfy the perturbative renormalization group equation:
    \[
    a(\beta) \sim \frac{1}{\Lambda} e^{-\frac{\beta}{2N\beta_0}}
    \]
\end{enumerate}

\begin{remark}[Scaling and Monotonicity]
In many physical arguments, one assumes that the string tension is monotonic in $\beta$. This would follow from correlation inequalities (like GKS inequalities) if they held for non-Abelian gauge theories. However, such inequalities have not been rigorously established for $SU(N)$ gauge theory. Thus, the monotonicity of the gap and string tension remains a conjecture.
\end{remark}
